% !TeX root = ../main.tex
% ===== 5 实验结果与分析 =====
\section{实验结果与分析}
\label{sec:r-results}

本文实验材料包括三张自然堆积真实图像及其到竞赛任务的端到端输出。真实图像没有配对的机械筛分逐筛称量、人工实例 mask 或材料形貌标签，因此本章报告场景输出、数量与质量代理统计差异、规则输出行为、软件一致性和计算耗时，不评估较真实结果的准确率。

\subsection{真实场景的级配输出}

本节以图\ref{fig:r-three-gsd}所示三组自然堆积样本为例，按 $D_{50}$ 从粗到细排列：粗骨料居多（图\ref{fig:r-three-gsd-a}）、中等骨料居多（图\ref{fig:r-three-gsd-b}）和细骨料居多（图\ref{fig:r-three-gsd-c}）三种代表性场景（对应表\ref{tab:r-demo-main}的样本 A、B、C），现有产物使用的尺度为 0.208~mm/px。三组图像共检测到 3858 个实例，并生成逐颗粒几何表、筛分等效粒径、质量代理级配、投影形貌、异常候选和场景级可视化。表\ref{tab:r-demo-main}给出异常规则筛选后的正常实例结果，三组数据统一使用 $\boldsymbol{\theta}=(1,0,0)$ 和 $\gamma=3$。

\begin{table}[t]
\centering
\caption{三张自然堆积样本的视觉质量代理级配输出（样本 A 为粗骨料居多、样本 B 为中等骨料居多、样本 C 为细骨料居多，分别对应图\ref{fig:r-three-gsd-a}--图\ref{fig:r-three-gsd-c}）。}
\label{tab:r-demo-main}
\footnotesize
\renewcommand{\arraystretch}{1.12}
\begin{tabular}{lrrrrll}
\toprule
样本 & 实例数 & 正常实例 & $D_{10}$ & $D_{50}$ & $D_{90}$ & 主粒级 \\
 & & & \multicolumn{3}{c}{/ mm} & \\
\midrule
样本 A（粗） & 968 & 703 & 15.6 & 27.5 & 38.0 & 25--31.5~mm（29.7\%） \\
样本 B（中等） & 1924 & 1193 & 11.1 & 23.6 & 30.6 & 25--31.5~mm（34.2\%） \\
样本 C（细） & 966 & 520 & 6.7 & 14.6 & 30.6 & 5--10~mm（32.7\%） \\
\bottomrule
\end{tabular}
\end{table}

按 $D_{50}$ 从粗到细排序，图\ref{fig:r-three-gsd-a}（样本 A，27.5~mm）最粗，图\ref{fig:r-three-gsd-b}（样本 B，23.6~mm）居中，图\ref{fig:r-three-gsd-c}（样本 C，14.6~mm）最细；图\ref{fig:r-three-gsd-a} 的 $D_{90}=38.0$~mm 也显示出更强的粗粒尾部。粒级组成与这一排序一致：图\ref{fig:r-three-gsd-a} 的 25--40~mm 约 55.6\%，图\ref{fig:r-three-gsd-b} 的 20--31.5~mm 合计约 55.3\%，而图\ref{fig:r-three-gsd-c} 的 5--16~mm 约 51.9\%。三组样本使用相同尺度口径、筛分映射和质量指数，因此这些差异反映了各场景的实例几何组成。

\begin{figure}[p]
\centering
\begin{subfigure}{0.95\linewidth}
\centering
\begin{minipage}[c]{0.66\linewidth}
\centering
\includegraphics[width=\linewidth]{agg_005_IG2_full_set_cp_SAM_pred_grains_re_scaled_gsd_comparison.png}
\end{minipage}\hfill
\begin{minipage}[c]{0.28\linewidth}
\centering
\includegraphics[width=\linewidth]{fig4_005_seg.jpg}
\end{minipage}
\caption{样本 A（粗骨料居多）}
\label{fig:r-three-gsd-a}
\end{subfigure}

\vspace{0.6em}

\begin{subfigure}{0.95\linewidth}
\centering
\begin{minipage}[c]{0.66\linewidth}
\centering
\includegraphics[width=\linewidth]{agg_001_IG2_full_set_cp_SAM_pred_grains_re_scaled_gsd_comparison.png}
\end{minipage}\hfill
\begin{minipage}[c]{0.28\linewidth}
\centering
\includegraphics[width=\linewidth]{fig4_001_seg.jpg}
\end{minipage}
\caption{样本 B（中等骨料居多）}
\label{fig:r-three-gsd-b}
\end{subfigure}

\vspace{0.6em}

\begin{subfigure}{0.95\linewidth}
\centering
\begin{minipage}[c]{0.66\linewidth}
\centering
\includegraphics[width=\linewidth]{agg_029_IG2_full_set_cp_SAM_pred_grains_re_scaled_gsd_comparison.png}
\end{minipage}\hfill
\begin{minipage}[c]{0.28\linewidth}
\centering
\includegraphics[width=\linewidth]{fig4_029_seg.jpg}
\end{minipage}
\caption{样本 C（细骨料居多）}
\label{fig:r-three-gsd-c}
\end{subfigure}
\caption{三组代表性场景的级配曲线（左）与实例分割叠加（右）对照，按 $D_{50}$ 从粗到细排列（样本 A 为粗骨料居多、样本 B 为中等骨料居多、样本 C 为细骨料居多，同表\ref{tab:r-demo-main}）。固定参数下，三组场景保持不同的分布中心和粗细尾部。}
\label{fig:r-three-gsd}
\end{figure}

\subsection{数量分布与质量代理分布}

为单独观察统计权重的影响，表\ref{tab:r-number-mass}使用每张图的\emph{全部检测实例}，不做异常剔除，只把实例权重从等权改为 $d^3$。三组样本的数量型 $D_{50}$ 为 5.2--7.6~mm，质量代理 $D_{50}$ 为 14.1--27.5~mm；对应增量分别为 17.5、19.9 和 8.9~mm。

\begin{table}[t]
\centering
\caption{同一全部实例集合上的数量型与质量代理型分位数比较（样本 A 为粗骨料居多、样本 B 为中等骨料居多、样本 C 为细骨料居多，同图\ref{fig:r-three-gsd}与表\ref{tab:r-demo-main}）。}
\label{tab:r-number-mass}
\footnotesize
\renewcommand{\arraystretch}{1.12}
\begin{tabular}{lrrrrrrr}
\toprule
& \multicolumn{3}{c}{数量型 / mm} & \multicolumn{3}{c}{质量代理 / mm} & $\Delta D_{50}$ \\
\cmidrule(lr){2-4}\cmidrule(lr){5-7}
样本 & $D_{10}$ & $D_{50}$ & $D_{90}$ & $D_{10}$ & $D_{50}$ & $D_{90}$ & / mm \\
\midrule
样本 A（粗） & 3.3 & 7.6 & 24.5 & 15.4 & 27.5 & 38.0 & 19.9 \\
样本 B（中等） & 3.1 & 6.0 & 17.3 & 10.6 & 23.5 & 30.6 & 17.5 \\
样本 C（细） & 3.0 & 5.2 & 9.2 & 5.9 & 14.1 & 30.6 & 8.9 \\
\bottomrule
\end{tabular}
\end{table}

这种整体向粗端的移动来自权重结构。自然堆积图像中，小颗粒往往在实例数量上占比高，但单颗粒质量贡献较低；粗颗粒数量较少，却会在 $d^3$ 权重下获得更高贡献。图\ref{fig:r-three-gsd-b}（样本 B）中，$<5$~mm 实例约占全部检测数量的 38.0\%；将这些目标从正常口径中排除后，质量代理 $D_{50}$ 仅由约 23.50~mm 变为 23.56~mm。大量小实例因此会显著改变数量分布，却未必以相同比例改变质量代理分布。

同一权重机制也会放大粗尺度 merge 的影响。两个相邻颗粒若被合并成一个偏大实例，偏大的 $d_i$ 还会通过 $d_i^3$ 获得更高统计权重，使误差进一步传递到粗端级配和高分位数。因此，上游实例边界质量与下游质量代理并非独立误差源，而是存在明确的放大关系。

这组对比说明数量分布与质量相关统计需要分开建模。$\gamma=3$ 在本文中仍是几何相似性先验；当前数据展示的是采用该先验后的统计行为，而不是对质量指数本身的机械称量校准结果。

\subsection{投影形貌与异常候选}

形貌和异常分支共享同一组实例几何。表\ref{tab:r-shape-anomaly}把两类辅助输出合并列出：普通投影形貌占 55.0\%--64.1\%，针片状候选占 14.0\%--20.9\%；异常规则则主要触发 $<5$~mm 小目标，三组样本的 $>50$~mm 和疑似泥团均为 0。

\begin{table}[t]
\centering
\caption{投影形貌比例与异常候选计数（样本 A 为粗骨料居多、样本 B 为中等骨料居多、样本 C 为细骨料居多，同图\ref{fig:r-three-gsd}；形貌四列为全部检测实例占比（\%），异常三列为实例数）。}
\label{tab:r-shape-anomaly}
\footnotesize
\renewcommand{\arraystretch}{1.10}
\begin{tabular}{lrrrrrrr}
\toprule
样本 & 针片状候选 & 圆形 & 棱角状 & 普通 & $<5$~mm & $>50$~mm & 疑似泥团 \\
\midrule
样本 A（粗） & 20.9 & 6.6 & 17.6 & 55.0 & 265 & 0 & 0 \\
样本 B（中等） & 16.3 & 8.0 & 16.4 & 59.3 & 731 & 0 & 0 \\
样本 C（细） & 14.0 & 11.0 & 11.0 & 64.1 & 446 & 0 & 0 \\
\bottomrule
\end{tabular}
\end{table}

\begin{figure}[H]
\centering
\begin{subfigure}{0.39\linewidth}\centering
\includegraphics[width=\linewidth]{agg_001_IG2_full_set_cp_SAM_pred_grains_re_scaled_shape_needle_flaky.png}
\caption{针片状候选}
\end{subfigure}\hspace{0.04\linewidth}%
\begin{subfigure}{0.39\linewidth}\centering
\includegraphics[width=\linewidth]{agg_001_IG2_full_set_cp_SAM_pred_grains_re_scaled_shape_round.png}
\caption{圆形}
\end{subfigure}

\vspace{0.4em}

\begin{subfigure}{0.39\linewidth}\centering
\includegraphics[width=\linewidth]{agg_001_IG2_full_set_cp_SAM_pred_grains_re_scaled_shape_angular.png}
\caption{棱角状}
\end{subfigure}\hspace{0.04\linewidth}%
\begin{subfigure}{0.39\linewidth}\centering
\includegraphics[width=\linewidth]{agg_001_IG2_full_set_cp_SAM_pred_grains_re_scaled_shape_regular.png}
\caption{普通}
\end{subfigure}
\caption{样本 B（图\ref{fig:r-three-gsd-b}）的投影形貌高亮。可视化用于核对规则命中的空间位置与轮廓特征。}
\label{fig:r-shape-montage}
\end{figure}

投影形貌比例用于描述当前视角下的二维轮廓组成，不对应三维针片状量规判定。异常分支在现有样本中能够检查 $<5$~mm 规则的执行，但缺少 $>50$~mm 或泥团阳性实例，因此这两类输出目前只具备规则级功能证据。

\subsection{软件一致性与运行效率}

当前项目实现了从等效粒径筛分到质量权重、加权分位数、形貌规则、异常优先级和报告生成的全链路闭合。

实例分割速度为 17~s、几何测量 9~s、筛分与报告统计小于 1~s，当前算法在部署阶段能实现约 27~s/张的推理速度。完整现场流程还包含尺度确认、图像质检、首次模型加载、结果复核以及必要的重拍或复跑。第\ref{sec:r-deploy}节据此安排现场操作预算。

综合以上结果，当前证据最直接支持两点：固定参数下，系统能够把不同自然堆积场景转换为可区分的逐颗粒与级配描述；在相同实例集合上，从数量等权切换到 $d^3$ 权重会系统性改变分布中心。标准机械筛分绝对误差、实例分割外部 P/R/F1 以及三维针片状或泥团材料识别准确率不在本文当前评价范围内。
